% vim: set textwidth=78 autoindent:

\section{Using external QGIS Python Plugins}\label{sec:external_plugins}\index{plugins}

% when the revision of a section has been finalized, 
% comment out the following line:
% \updatedisclaimer

External QGIS plugins are written in python. They are stored in 
an official, moderated repository and maintained by the individual 
author. Table \ref{tab:external_plugins} shows a list of plugins 
currently available with a short description.\footnote{Also updates of 
core plugins may be available in this repository as external overlays.} 
\footnote{The Python Plugin Installer is also an external Python Plugin, but 
it is part of the QGIS sources and automatically loaded and selectable inside
the QGIS Plugin Manager (see Section~\ref{sec:load_external_plugin}).}

When this manual was released, the external moderated QGIS plugin repository
was not fully established. A detailed documentation about the usage, the
author and other important information are provided with the external plugin
itself and is not part of this manual.

You find an up-to-date list of moderated external plugins in the QGIS
Official Repository of the \dropmenuopttwo{plugin_installer}{Fetch Python
Plugins...}) and at \url{http://qgis.osgeo.org/download/plugins.html}. 

\begin{table}[H]
\centering
\caption{Current moderated external QGIS Plugins}\label{tab:external_plugins}\medskip
\small
 \begin{tabular}{|l|l|p{4in}|}
\hline \textbf{Icon} & \textbf{external Plugin} & \textbf{Description}\\
\hline
\includegraphics[width=0.7cm]{zoom2point_icon}
 & Zoom To Point \index{plugins!Zoom To Point} & Zooms to a coordinate 
  specified in the input dialog. You can specify the zoom level as well to 
  control the view extent.\\
\hline
\end{tabular}
\end{table}

A detailed install description for external python plugins can be found in 
Section \ref{sec:load_external_plugin}.

\minisec{User-Contributed Python Plugin Repository and author repositories}

Apart from the moderated external plugins there exists another unofficial
Python Plugin repository. It contains plugins that are not yet mature
enough to include them to the official repository, however some of them
can be quite useful. Furthermore a few of our contributors maintain
their own repositories.

To add the unofficial repository and the author repositories open the 
Plugin Installer (\mainmenuopt{Plugins} > \dropmenuopttwo{plugin_installer}{Fetch Python Plugins...}),
go to the \tab{Repositories} tab and click the \button{Add 3rd party repositories}
button. If you don't want one or more of the added repositories, disable it with the
\button{Edit...} button or completely remove with the \button{Delete} button.

\begin{Tip} \caption{\textsc{Add more external Plugins}}
\qgistip{Apart from the official QGIS plugin repository you can add 
more external repositories. Therefore select the Repositories tab in 
the Python Plugins Installer
}
\end{Tip}

